\section{Related Work}
\label{sec:related}

Our contribution sits at the intersection of three literatures that have
developed largely independently: distribution-free uncertainty quantification
via conformal prediction, fuzzy neural architectures for graphs, and machine
learning for network security. We review each in turn, and in each case we
identify not merely what has been done but what the accumulated body of work
leaves unresolved for the deployment setting we target.

\subsection{Conformal Prediction — Foundations}

Conformal prediction (CP) \citep{angelopoulos21} converts any point predictor
into a set-valued predictor with a finite-sample, distribution-free coverage
guarantee: for a user-specified error rate $\alpha$, the returned set satisfies
$\Prob(y \in \Calpha(x)) \ge 1-\alpha$ irrespective of the underlying data
distribution and of whether the base model is well specified. The guarantee
rests on a single assumption — exchangeability between the calibration sample
and the test point — and this is precisely what makes CP attractive for
security applications, where model misspecification is the norm rather than the
exception and where distributional assumptions about attacker behaviour are
untenable.

The exchangeability requirement is also CP's principal limitation, and a
substantial line of work has sought to relax it. \citet{barber23} develop
weighted conformal prediction, which restores approximate validity under
covariate shift by reweighting calibration points according to a likelihood
ratio between training and test distributions; the coverage loss is bounded by
the total-variation distance between the two. This matters here because
security traffic is non-stationary by construction: adversaries adapt, network
topologies evolve, and the deployment distribution drifts away from whatever
was captured at calibration time. \citet{cp_survey} survey the field from a
data-centric perspective and note that graph-structured data poses a distinctive
difficulty, since the interdependence induced by edges violates the i.i.d.\
intuition that motivates the exchangeability condition even when the condition
itself can be salvaged.

Within cybersecurity specifically, the present authors have applied CP to
intrusion detection \citep{dang23ids} and botnet detection
\citep{dang24botnet}, and developed CONFIDE for causal competing-risk settings
\citep{dang26confide}. Those works operate on tabular flow representations and
therefore inherit neither the benefits nor the complications of graph structure;
the present paper is a deliberate step towards the graph setting, where the
relational structure that makes attacks detectable is also what complicates the
statistical guarantee.

\subsection{Conformal Prediction for Graph Neural Networks}

Transferring CP to graphs is not a matter of applying the standard recipe to a
new model class. In transductive node classification, the prediction for one
node depends on the features and labels of others through message passing, so
calibration and test scores are not independent, and the exchangeability
argument must be re-established rather than assumed.

\citet{cfgnn} resolved this with CF-GNN, now the de-facto baseline. Their key
technical contribution is a \emph{permutation-invariance condition}: if the
model treats calibration and test nodes symmetrically — which holds for standard
message-passing architectures under a random calibration/test split — then the
nonconformity scores are exchangeable and split CP remains valid despite the
statistical coupling induced by the graph. They further propose a topology-aware
correction that shrinks prediction sets by exploiting the observation that a
node's optimal set size correlates with its neighbourhood's.

Subsequent work has largely pursued efficiency, taking validity as settled.
\citet{daps} propose DAPS, which diffuses conformity scores across edges under a
homophily assumption, on the intuition that adjacent nodes should receive
similar scores; this tightens sets when homophily holds. \citet{snaps} pursue a
related idea from a different angle, aggregating scores from nodes that are
similar in \emph{feature} space rather than adjacent in the graph, improving the
singleton-hit proportion — the fraction of nodes receiving an unambiguous
single-class prediction, which is what an analyst can act on directly.
\citet{clarkson23} addresses the inductive regime, where test nodes arrive after
calibration and the transductive exchangeability argument no longer applies
directly.

Our own results (Section~\ref{sec:experiments}) suggest a caveat that this line
of work has not emphasised: score-smoothing operations that improve average
efficiency can degrade calibration, because the smoothed score is no longer a
function of the node alone and the exchangeability argument that licensed the
guarantee is weakened. We observe exactly this behaviour for the
diffusion-based and similarity-based methods under chronological splits.

Closest to our work is RR-GNN \citep{rrgnn}, which introduces
\emph{Graph-Structured Mondrian CP}. Rather than computing one global quantile,
RR-GNN partitions nodes into graph communities and computes a quantile per
community, yielding coverage conditional on community membership. This is the
right structural insight for our setting, and FMCP builds directly on it: a
marginal guarantee averaged over all traffic can be satisfied while a particular
attack family is systematically under-covered, and it is precisely the rare
families that matter operationally. RR-GNN, however, uses a residual-reweighted
softmax score, incorporates no fuzzy reasoning, produces no human-readable
explanation, and is not evaluated on cybersecurity data.

Further work extends CP-GNNs to federated settings \citep{akgul25}, temporal
graphs \citep{ncp2025,ncpnet25}, and conformal-aware training
\citep{confrank25}, in which the CP objective is folded into the loss so the
model learns to produce efficiently calibratable scores rather than being
calibrated post hoc. \citet{crcs_gad} extend conformal risk control to graph
anomaly detection, targeting a bound on a general loss rather than on
miscoverage. \citet{fuzzy_cp_evalue} construct \emph{fuzzy} prediction sets via
e-values in the non-graph setting, replacing the binary in/out decision with a
graded membership; this is the closest theoretical antecedent to combining fuzzy
reasoning with conformal inference, although it does not address graphs and does
not use a fuzzy integral as the nonconformity measure.

\subsection{Fuzzy Graph Neural Networks}

Fuzzy set theory offers a principled treatment of the graded, partially
overlapping category membership that characterises network telemetry, where the
boundary between benign and malicious traffic is genuinely indistinct rather
than merely uncertain.

\citet{flgnn} proposed FL-GNN, which embeds a Cartesian-product fuzzy rule base
inside the message-passing layer so that aggregation is performed by rule firing
rather than by a fixed permutation-invariant operator. The appeal is that the
learned rules are inspectable: each is an IF--THEN statement over membership
functions, so a prediction can be traced to the rules that produced it. Our
fuzzy message-passing layer extends this construction with Gaussian membership
functions, the product T-norm, and centroid defuzzification.

Fuzzy reasoning has also been combined with attention. \citet{fgat} use fuzzy
rough-set attention weights with dynamic negative sampling for link prediction,
and \citet{mfgat} extend this to multiple views with learnable global pooling.
\citet{firegnn} pursue a neuro-symbolic direction with trainable IF--THEN rules,
and \citet{fgenconv} apply explainable fuzzy GNNs to leak detection in water
distribution networks — an infrastructure-anomaly task structurally analogous to
ours, in that both require an operator to act on an alert rather than merely
record it. The present authors applied fuzzy embeddings to software-defined
network intrusion detection in \citep{dang24fuzzy}.

The survey of \citet{fgnn_survey} is directly relevant to our motivation. It
identifies the absence of formal guarantees as an open problem cutting across
the entire fuzzy GNN literature: these architectures quantify uncertainty in the
sense of producing graded memberships, but a membership value carries no
frequentist interpretation and supports no statement about how often the true
class is captured. Fuzziness and calibration are distinct properties, and
providing the former does not supply the latter. This is the gap FC-GNN
addresses, and it is why the conformal layer is not an optional add-on but the
component that converts graded output into an actionable guarantee.

\subsection{GNNs for Cybersecurity}

Graph representations are natural for network security because attacks manifest
relationally: lateral movement, command-and-control beaconing and distributed
denial of service are all defined by communication structure rather than by any
single flow.

\citet{graphids} report 99.98\% PR-AUC on NF-BoT-IoT with a self-supervised
E-GraphSAGE encoder combined with a transformer autoencoder, training without
attack labels — an important property given how quickly labelled attack data
becomes stale. \citet{gnnids} combine a static attack graph encoding known
vulnerability paths with dynamic telemetry, and \citet{xgnid} pair a
heterogeneous GNN with a language model to generate natural-language alert
descriptions. The survey of \citet{gnn_ids_survey} covers the broader landscape.
The present authors previously proposed GNN-based NIDS baselines
\citep{dang23wifi} and federated IoT malware detection \citep{dang24federated}.

A recurring pattern in this literature is that headline accuracy is very high —
frequently above 99\% — while deployment experience reports persistent
false-alarm burden. The discrepancy is informative: aggregate accuracy on a
benchmark with a dominant benign class says little about behaviour on the rare
classes that generate actionable alerts, which is an argument for reporting
class-conditional or community-conditional quantities rather than aggregates,
and is the motivation for the conditional coverage gap we adopt as our primary
calibration metric.

\rev{As GNNs see wider deployment in cybersecurity, their robustness against
graph adversarial threats has become a critical concern. Recent literature
explores vulnerabilities such as node injection attacks and corresponding
defences (e.g., GZOO, DShield) \citep{gzoo2024, dshield2023}. Node injection is
the realistic threat model for network telemetry, since an attacker controls the
hosts it owns and the flows it emits but not the records of other hosts; such an
attack bears on FC-GNN at two points, in that it violates the exchangeability
premise of the conformal guarantee and perturbs the topology on which the
community partition depends. Furthermore, structural entropy has emerged as a
complementary theoretical perspective for understanding graph complexity and
structural anomalies, which is highly relevant for proactive botnet detection
\citep{zeng2025structural}. This is of direct interest for the taxon-selection
problem that Mondrian CP leaves open: modularity-based partitioning optimises at
a single resolution, whereas minimising structural entropy yields a hierarchy in
which granularity becomes an explicit, tunable level.}

\subsection{Conformal Prediction in Cybersecurity}

A smaller body of work applies CP directly to security tasks.
\citet{botnet_cp} add a conformal rejection layer to a botnet classifier,
abstaining when the prediction set is uninformative and thereby rejecting 58\%
of would-be incorrect predictions on ISCX — a concrete demonstration that
calibrated abstention is more useful operationally than a forced decision.
\citet{damba26} extend open-set CP to zero-day detection in industrial IoT,
where the true class may be absent from the training taxonomy altogether.
\citet{cp_cps26} propose sliding-window temporal quantile CP for cyber-physical
systems, refreshing the conformal threshold as the process evolves.

\citet{cp_drift} report the most important negative result for our purposes:
under heavy concept drift, CP-derived pseudo-labels become inconsistent, so the
guarantee degrades precisely in the regime where it is most needed. This
directly motivates our chronological split protocol, under which calibration and
test data are separated in time rather than sampled at random, and it is the
reason we regard random-split evaluations of CP-GNNs as optimistic for security
deployment.

\subsection{Fuzzy Methods in Cybersecurity}

Fuzzy reasoning has a long history in intrusion detection, predating the current
generation of graph methods. \citet{fuzzy_iot} combine fuzzy logic with CNN
ensembles for IoT NIDS, reporting 99.72\% on NSL-KDD; \citet{fuzzy_sdn} optimise
a fuzzy-rule deep architecture for SDN-enabled IoT using a hunger-games search;
and \citet{gfs_gan} evolve fuzzy rules under adversarial training, using
generated attack traffic to harden the rule base.

These results establish that fuzzy reasoning is practically effective for
security data and that its rule bases are valued for interpretability. None
provides a coverage guarantee. Taken together with the CP literature reviewed
above — which provides guarantees but produces opaque scores and no rule-level
explanation — this defines the gap the present work addresses: the two
properties have been pursued separately, and an operator needs both from the
same system.
